\documentclass[../../../main.tex]{subfiles}

\begin{document}

\section{Limits and Continuity}

Before we get into limits and continuity of multivariable functions,
let's firsts review what limits are for single variable functions.
Please note that the contents here will not be as ``formal'' or
rigorous as taught in Real Analysis courses, but maybe I could
include them in the future! First, recall the following definition
of limit for a single variable function.

\begin{definition}{}{}
Given a function $f(x)$, the \textbf{limit} of $f(x)$ as $x$
approaches $a$ denoted as $\lim_{x \to a} f(x) = L$ is defined
as the following:
\vspace{0.8\baselineskip}
\begin{center}
\begin{minipage}{0.7\textwidth}
Given any $\epsilon > 0$, there exists $\delta > 0$ such that
if $0 < | x - a | < \delta$, then $| f(x) - L | < \epsilon$.
\end{minipage}
\end{center}
\end{definition}

An intuitive way of understanding this epsilon-delta definition of
limit would be finding a value that the bound of $f(x)$ gets
arbitrarily close to as the distance between the bound for $x$ gets
arbitrarily close to zero. Now for $2$-space, you are probably
familiar with a graph involving a shaded rectangle. For $3$-space,
it is similar, be we can think of it as a circle. Our bound is no
longer the width of the rectangle, but the radius of the circle
surrounding the point. Consider the following definition for
functions in $3$-space, or functions with two variables.

\begin{definition}{}{}
Given a function $f(x, y)$, the limit of $f(x, y)$ as $(x, y)$
approaches $(x_0, y_0)$ denoted as $\lim_{(x, y) \to (x_0, y_0)}
f(x, y) = L$ is defined as the following:
\vspace{0.6\baselineskip}
\begin{center}
\begin{minipage}{0.8\textwidth}
Given any $\epsilon > 0$, there exists $\delta > 0$ such that
for all $(x, y) \in \dom(f)$, if $0 < \sqrt{(x - x_0)^2 + (y - y_0)^2}
< \delta$, then $| f(x, y) - L | < \epsilon$.
\end{minipage}
\end{center}
\end{definition}

Considering that $|x - a|$ is the distance between $x$ and $a$,
I hope the expression $\sqrt{(x - x_0)^2 + (y - y_0)^2}$ is a
natural extension because it represents the distance between
$(x, y)$ and $(x_0, y_0)$. Let's take at an example.

\begin{problem}{}{}
Show that $\lim_{(x, y) \to (0, 0)} \frac{x + y}{x^2 + 1} = 0$.
\end{problem}
\begin{solution}
Given $\epsilon > 0$, choose $\delta = \frac{\epsilon}{\sqrt{2}}$.
Then, the following equations hold by AM-GM inequality.
\[
\left| \frac{x + y}{x^2 + 1} - 0 \right|
\leq | x + y |
\leq \sqrt{2x^2 + 2y^2}
< \sqrt{2} \delta
= \epsilon
\]
Therefore, $0 < \sqrt{(x - x_0)^2 + (y - y_0)^2} < \delta$,
then $| f(x, y) - 0 | < \epsilon$.
\end{solution}

Similarly, we can replace the points in $2$-space to $3$-space to
generalize for function with three variables. Another way of thinking
about limits for multivariable functions are through different curves.

\begin{definition}{}{}
Given a function $f(x, y)$ and a smooth curve $C$ in $2$-space with
$x(t) = x$ and $y(t) = y$ and $(x_0, y_0) = (x(t_0), y(t_0))$,
then the following equation hold.
\[
\lim_{(x, y) \underset{\to}{C} (x_0, y_0)} f(x, y)
= \lim_{t \to t_0} f(x(t), y(t))
\]
\end{definition}

This is now easier because $\lim_{t \to t_0} f(x(t), y(t))$ is now a
single variable limit. Here is an example continuing with our previous
one.

\begin{problem}{}{}
Compute $\lim_{(x, y) \to (0, 0)} \frac{x + y}{x^2 + 1}$ along the
path $y = x$.
\end{problem}
\begin{solution}
Notice that $y = x$ can be represented as $x(t) = t$ and $y(t) = t$.
Therefore, the limit can be rewritten as $\lim_{t \to 0}
\frac{t + t}{t^2 + 1} = 0$. Therefore, the limit along the path
$y = x$ is zero.
\end{solution}

However an important thing to note is that we have infinitely many
different ways to represent a limit along different curves because
there are infinitely different curves. This can cause a problem.
What if the limits return different values from one to another?

\begin{important}
Finding limits along a curve does not represent the function's limit
and cannot be used to prove the limit.
\end{important}

This is important because it shows that we can use this method to
disprove the existence of limit, but not find limit. Here is an
example.

\begin{problem}{}{}
Find, if any, the value of $\lim_{(x, y) \to (0, 0)}
\frac{3y}{x^2 + y}$.
\end{problem}
\begin{solution}
First, consider the limit along a curve $y = kx^2$ for any $k$.
Notice that the limit can be represented as $\lim_{t \to 0}
\frac{3kx^2}{x^2 + kx^2} = \frac{3k}{1 + k}$. Because
$\frac{3k}{1 + k}$ depends on the value of $k$ and is not constant,
the limit does not exist.
\end{solution}

One final thing before we wrap up our discussions about limit is that
the properties of sum, different, product, and quotient applies the
same for the generalized limits and limits along different curves
just as the limit for single variable function!






\end{document}


